% v11 RQ3: design factors and calibration (AB1-AB5, H1, M6, DI2).
\section{RQ3: Design Factors and Calibration}\label{sec:design-results}

\subsection{The Default Does Not Dominate Its Predecessor}
On real edit-targeted public histories, the earlier cascade v1 relocates more targets than the final default (\smPubTgtVoneAccPct\% versus \smPubTgtCasAccPct\%; internal \smIntTgtVoneAccPct\% versus \smIntTgtCasAccPct\%), although the final default leads on controlled edits (\smCtlCasAccPct\% versus \smCtlVoneAccPct\%). The reversal is a primary result, not a footnote: the default threshold of 0.85 was tuned on controlled token edits and is too strict for naturally edited statements, on which it relocates only \thetaEditedEightyFive\% of edited statements correctly (Table~\ref{tab:tradeoff}). v1 also differs in tier order, search scope and threshold, so the comparison does not isolate any single choice.

\subsection{Factorial Ablation}
To separate the choices, we evaluated every combination of tier order (local fuzzy search before cross-file exact search, as deployed, or after it, as in v1), search scope (a window around the diff-predicted line or the whole file) and threshold $\theta\in\{0.50,0.55,\dots,0.95\}$ on both controlled seeds and all real-history observations. The two corners matching the deployed cascade and v1 reproduce the archived outcomes exactly. Averaged over thresholds, searching the whole file instead of the local window lowers the rate at which deleted code is correctly reported as removed by \wpaScopeGoneInt{} points on internal and \wpaScopeGonePub{} points on public projects: a distant look-alike is accepted in place of an honest \lbl{removed}. Running cross-file exact search before local fuzzy search raises missed re-checks by \wpaOrderMissInt{} and \wpaOrderMissPub{} points, because an edited block is then matched to an identical copy elsewhere, which is exactly the failure analyzed in Section~\ref{sec:safety}. Both deployed choices are therefore supported on safety grounds. Accuracy alone favours lower thresholds, and at equal thresholds the final design leads v1 in all nine projects after calibration; on internal projects the equal-threshold advantage of \wpaEqThrDiff{} points has a project-cluster interval of $[\wpaEqThrLo, \wpaEqThrHi]$, so its direction is not established there. Table~\ref{tab:factorial} reports the averaged contrasts, and the evidence package retains every cell; these are matched contrasts on shared observations, not randomized experiments.

\subsection{The Relocation--Deletion Trade-off}
Lowering the threshold recovers more edited statements and reports fewer deleted ones correctly (Table~\ref{tab:tradeoff}). At $\theta=0.50$, \thetaEditedFifty\% of edited statements are relocated but only \thetaGoneFifty\% of deleted statements are reported removed; at the default 0.85 the figures are \thetaEditedEightyFive\% and \thetaGoneEightyFive\%. No threshold is best for both. A deleted statement that is relocated attaches a concern to different code, which is the more consequential failure, whereas an edited statement that is not relocated becomes a visible \lbl{removed} anchor and requests re-check. This trade-off is intrinsic to similarity thresholding and is the central practical decision when deploying the method.

\subsection{Cross-Project Calibration}
In leave-one-project-out calibration over the nine projects, choosing the threshold that maximizes accuracy on eight projects selects $\theta=\wpaLopoTheta$ in all \wpaLopoFolds{} folds, both for the threshold alone and for the full factorial cell, and for both training objectives (Table~\ref{tab:lopo}). The selection is stable across projects. Applied to public edit-targeted histories, $\theta=\wpaLopoTheta$ relocates \smPubTgtLopoAccK{} of \smPubTgtLopoAccN{} targets (\smPubTgtLopoAccPct\%), against \smPubTgtCalAccPct\% at the archived calibrated value of 0.50 and \smPubTgtCasAccPct\% at the default. These histories had been inspected in earlier revisions, so this is retrospective transfer evidence, not a prospective held-out result.

\subsection{Safety-Constrained Calibration Did Not Transfer}
Calibrating for accuracy is misaligned with the stated objective, so we also chose thresholds on the four internal projects to maximize accuracy subject to a missed re-check rate of at most $\alpha$ and a deleted-reported-present rate of at most $\beta$, for $\alpha\in\{1,2,5\}\%$ and $\beta\in\{5,10,20\}\%$. Of the design and limit combinations, \wpaConstrFeasible{} were feasible on the internal projects and \wpaConstrInfeasible{} were not. Evaluated once on the public projects, \wpaConstrViolate{} of the \wpaConstrFeasible{} feasible choices violated at least one of their own limits. For example, $\theta=\wpaConstrTheta$ met an internal 20\% cap on deleted code reported present but reached \wpaConstrPubGonePresent\% on public edit-targeted observations. Empirical rates from one group of projects therefore do not bound the rates in another; a deployment that needs a guaranteed limit would require a calibration method with explicit transfer assumptions, and the projects here do not justify exchangeability.

\subsection{Choosing a Threshold in Practice}
For an institution with no history of its own, the evidence supports three steps. Keep the local search scope and the deployed tier order, which are favoured on safety grounds in every stratum. Treat $\theta$ as a policy choice between relocating edited code and reporting deleted code honestly: $\theta\approx\wpaLopoTheta$ maximized accuracy in every held-out project, whereas the default 0.85 favours honest removal. Monitor the outcomes after deployment, since supervisors' own corrections of misplaced anchors are the data a local calibration would need.
